\documentclass[12pt,a4paper]{article}
\usepackage[utf8]{inputenc}
\usepackage[T1]{fontenc}
\usepackage{geometry}
\geometry{a4paper, margin=1in}
\usepackage{hyperref}
\usepackage{graphicx}
\usepackage{booktabs}
\usepackage{longtable}
\usepackage{listings}
\usepackage{xcolor}
\usepackage{titlesec}
\usepackage{enumitem}

% Setup listings for code blocks
\definecolor{lightgray}{rgb}{0.95,0.95,0.95}
\lstset{
    backgroundcolor=\color{lightgray},
    basicstyle=\ttfamily\small,
    breaklines=true,
    frame=single,
    tabsize=2,
    showstringspaces=false
}

\title{\textbf{Self-Hosted Shorebird Architecture} \\ \Large Comprehensive Backend Specification \& Integration Guide}
\author{Generated by Shorebird AI Architect}
\date{\today}

\begin{document}

\maketitle
\tableofcontents
\newpage

\section{Introduction}
\subsection{Overview}
Shorebird is a revolutionary Code Push solution for Flutter applications, allowing developers to deploy over-the-air (OTA) updates seamlessly. Currently, Shorebird operates entirely as a managed cloud service. However, enterprise requirements such as data residency, strict compliance, and air-gapped environments necessitate a self-hosted, on-premise implementation.

This document details the complete architectural blueprint, API specifications, database schemas, and integration strategy required to construct a self-hosted backend that perfectly mimics the \texttt{shorebird\_code\_push\_protocol}. By implementing this specification, the modified Shorebird CLI and the native device Updater libraries will operate against your private infrastructure without requiring underlying structural changes to their core logic.

\subsection{Goals and Scope}
The primary objective is to replace \texttt{api.shorebird.dev} and \texttt{download.flutter.io} (via Shorebird's storage buckets) with your own privately managed infrastructure.
\begin{itemize}
    \item \textbf{Custom API Server:} A backend service handling the Shorebird Code Push Protocol.
    \item \textbf{Object Storage:} An S3-compatible storage layer (e.g., MinIO or AWS S3) to host engine binaries, release artifacts, and patch files.
    \item \textbf{Relational Database:} A persistent store for metadata (users, apps, releases, patches).
    \item \textbf{Artifact Proxy:} A configured proxy to intercept Flutter SDK requests and serve self-hosted engine artifacts.
\end{itemize}

\newpage

\section{High-Level Architecture}
The self-hosted Shorebird ecosystem consists of four main actors: the Developer (CLI), the API Server (Backend), the Storage Layer, and the End User (Mobile App/Updater).

\subsection{Component Diagram \& Data Flow}
\begin{enumerate}
    \item \textbf{Developer CLI (\texttt{shorebird\_cli}):} Configured with the \texttt{SHOREBIRD\_HOSTED\_URL} environment variable. It communicates with the custom API to create apps, publish releases, and upload patches.
    \item \textbf{Custom API Server:} Exposes the REST API defined by the \texttt{shorebird\_code\_push\_protocol}. It validates requests, manages the database, and generates presigned URLs for secure artifact uploads/downloads.
    \item \textbf{Relational Database (PostgreSQL):} Stores relationships between Organizations, Users, Apps, Releases, Patches, and Channels.
    \item \textbf{Object Storage (MinIO / S3):} The ultimate destination for binary artifacts. The \texttt{artifact\_proxy} service points here via \texttt{FLUTTER\_STORAGE\_BASE\_URL}.
    \item \textbf{Mobile App (Updater Library):} Periodically polls the custom API Server's \texttt{/api/v1/patches/check} endpoint to determine if a new patch is available for its specific release version and architecture.
\end{enumerate}

\newpage

\section{Backend Project Structure}
To implement the Custom API Server, a modern, statically typed language like Go, Node.js (TypeScript), or Rust is recommended. Below is the proposed directory structure for a scalable Go-based backend implementation.

\begin{lstlisting}[language=bash, caption={Recommended Backend Project Structure}]
shorebird-self-hosted-backend/
├── cmd/
│   └── server/
│       └── main.go              # Application entry point
├── internal/
│   ├── api/
│   │   ├── handlers/            # HTTP handlers (e.g., app_handler.go, patch_handler.go)
│   │   ├── middleware/          # Auth, CORS, Logging
│   │   └── router.go            # Route definitions
│   ├── config/                  # Environment and config loading
│   ├── database/
│   │   ├── migrations/          # SQL migration files
│   │   └── repository/          # Database interaction layer
│   ├── models/                  # Structs matching shorebird_code_push_protocol
│   ├── service/                 # Core business logic (Release matching, patching)
│   └── storage/                 # S3/MinIO interaction (Presigned URLs)
├── pkg/
│   └── utils/                   # Helper functions (Hashing, validation)
├── docker-compose.yml           # Local dev setup (Postgres + MinIO + API)
├── Dockerfile                   # Backend container definition
└── go.mod
\end{lstlisting}

\newpage

\section{Database Schema Design (PostgreSQL)}
The backend requires a robust relational schema to track the state of releases and patches across different platforms and tracks (channels).

\subsection{Table: \texttt{users}}
Stores authentication and profile data.
\begin{itemize}
    \item \texttt{id} (UUID, Primary Key)
    \item \texttt{email} (VARCHAR, Unique)
    \item \texttt{display\_name} (VARCHAR)
    \item \texttt{created\_at} (TIMESTAMP)
\end{itemize}

\subsection{Table: \texttt{organizations}}
Groups apps and manages billing/access in the public version; useful for multi-tenant self-hosting.
\begin{itemize}
    \item \texttt{id} (UUID, Primary Key)
    \item \texttt{name} (VARCHAR)
\end{itemize}

\subsection{Table: \texttt{apps}}
Represents a Flutter project.
\begin{itemize}
    \item \texttt{id} (VARCHAR, Primary Key) - Matches the Shorebird App ID generated by the CLI.
    \item \texttt{org\_id} (UUID, Foreign Key)
    \item \texttt{display\_name} (VARCHAR)
\end{itemize}

\subsection{Table: \texttt{releases}}
A specific, immutable build of an app (e.g., version 1.0.0+1).
\begin{itemize}
    \item \texttt{id} (SERIAL, Primary Key)
    \item \texttt{app\_id} (VARCHAR, Foreign Key)
    \item \texttt{version} (VARCHAR) - e.g., "1.0.0+1"
    \item \texttt{flutter\_revision} (VARCHAR) - The engine commit hash.
    \item \texttt{display\_name} (VARCHAR)
    \item \texttt{created\_at} (TIMESTAMP)
\end{itemize}

\subsection{Table: \texttt{release\_artifacts}}
The binaries (AAR, IPA, APK, AppBundle) associated with a release.
\begin{itemize}
    \item \texttt{id} (SERIAL, Primary Key)
    \item \texttt{release\_id} (INT, Foreign Key)
    \item \texttt{arch} (VARCHAR) - e.g., "aarch64", "x86\_64"
    \item \texttt{platform} (VARCHAR) - e.g., "android", "ios"
    \item \texttt{hash} (VARCHAR) - SHA256 of the artifact.
    \item \texttt{size} (INT)
    \item \texttt{url} (VARCHAR) - S3 Object Key or URL.
\end{itemize}

\subsection{Table: \texttt{patches}}
A delta update applied to a specific release.
\begin{itemize}
    \item \texttt{id} (SERIAL, Primary Key)
    \item \texttt{app\_id} (VARCHAR, Foreign Key)
    \item \texttt{number} (INT) - Sequential patch number (1, 2, 3...).
\end{itemize}

\subsection{Table: \texttt{patch\_artifacts}}
The actual binary patch files applied to the device.
\begin{itemize}
    \item \texttt{id} (SERIAL, Primary Key)
    \item \texttt{patch\_id} (INT, Foreign Key)
    \item \texttt{arch} (VARCHAR)
    \item \texttt{platform} (VARCHAR)
    \item \texttt{hash} (VARCHAR)
    \item \texttt{size} (INT)
    \item \texttt{url} (VARCHAR)
\end{itemize}

\subsection{Table: \texttt{channels}}
Deployment tracks (e.g., "stable", "beta", "staging").
\begin{itemize}
    \item \texttt{id} (SERIAL, Primary Key)
    \item \texttt{app\_id} (VARCHAR, Foreign Key)
    \item \texttt{name} (VARCHAR)
\end{itemize}

\subsection{Table: \texttt{release\_patches}}
Mapping a patch to a release and a specific channel.
\begin{itemize}
    \item \texttt{release\_id} (INT, Foreign Key)
    \item \texttt{patch\_id} (INT, Foreign Key)
    \item \texttt{channel\_id} (INT, Foreign Key)
    \item \textbf{Primary Key:} (release\_id, patch\_id, channel\_id)
\end{itemize}

\newpage

\section{REST API Specification}
The custom backend must implement the following endpoints to satisfy the \texttt{shorebird\_code\_push\_client}. All endpoints are prefixed with \texttt{/api/v1}. Authentication is handled via Bearer Tokens (JWT).

\subsection{App Management}

\subsubsection{Create App}
\textbf{Endpoint:} \texttt{POST /api/v1/apps} \\
\textbf{Description:} Registers a new application. \\
\textbf{Request Body:}
\begin{lstlisting}[language=json]
{
  "display_name": "My Flutter App"
}
\end{lstlisting}
\textbf{Response (200 OK):}
\begin{lstlisting}[language=json]
{
  "id": "app-uuid-string",
  "display_name": "My Flutter App"
}
\end{lstlisting}

\subsection{Release Management}

\subsubsection{Create Release}
\textbf{Endpoint:} \texttt{POST /api/v1/apps/:appId/releases} \\
\textbf{Description:} Creates a new base release for the app. \\
\textbf{Request Body:}
\begin{lstlisting}[language=json]
{
  "version": "1.0.0+1",
  "flutter_revision": "856815ea3fb...",
  "display_name": "v1.0.0"
}
\end{lstlisting}
\textbf{Response (200 OK):} Returns the Release object.

\subsubsection{Create Release Artifact}
\textbf{Endpoint:} \texttt{POST /api/v1/releases/:releaseId/artifacts} \\
\textbf{Description:} Registers a binary for a release and returns a presigned URL for the CLI to upload the file to S3. \\
\textbf{Request Body:}
\begin{lstlisting}[language=json]
{
  "arch": "aarch64",
  "platform": "android",
  "hash": "sha256-hash-string",
  "size": 15000000
}
\end{lstlisting}
\textbf{Response (200 OK):}
\begin{lstlisting}[language=json]
{
  "id": 1,
  "release_id": 123,
  "arch": "aarch64",
  "platform": "android",
  "hash": "sha256-hash-string",
  "size": 15000000,
  "url": "https://s3.yourdomain.com/presigned-upload-url-xyz"
}
\end{lstlisting}

\subsection{Patch Management}

\subsubsection{Create Patch}
\textbf{Endpoint:} \texttt{POST /api/v1/apps/:appId/patches} \\
\textbf{Description:} Initializes a new patch record. \\
\textbf{Request Body:} None required, or optional metadata. \\
\textbf{Response (200 OK):} Returns a Patch object with the next sequential \texttt{number}.

\subsubsection{Create Patch Artifact}
\textbf{Endpoint:} \texttt{POST /api/v1/patches/:patchId/artifacts} \\
\textbf{Description:} Registers a patch binary and provides an upload URL. \\
\textbf{Request Body:} Similar to Release Artifact. \\
\textbf{Response (200 OK):} Returns Patch Artifact with S3 presigned upload URL.

\subsubsection{Promote Patch}
\textbf{Endpoint:} \texttt{POST /api/v1/patches/:patchId/promote} \\
\textbf{Description:} Assigns a patch to a specific release and channel, making it active. \\
\textbf{Request Body:}
\begin{lstlisting}[language=json]
{
  "channel_id": 1,
  "release_id": 123
}
\end{lstlisting}
\textbf{Response (200 OK):} Returns success.

\subsection{Updater Endpoint (Critical)}

\subsubsection{Check for Patches}
\textbf{Endpoint:} \texttt{POST /api/v1/patches/check} \\
\textbf{Description:} This endpoint is called directly by the mobile app's Updater library (written in Rust). It checks if a newer patch is available for the device. \\
\textbf{Request Body:}
\begin{lstlisting}[language=json]
{
  "app_id": "app-uuid-string",
  "channel": "stable",
  "release_version": "1.0.0+1",
  "patch_number": 0,
  "arch": "aarch64",
  "platform": "android"
}
\end{lstlisting}
\textbf{Response (200 OK):}
\begin{lstlisting}[language=json]
{
  "patch_available": true,
  "patch": {
    "number": 1,
    "download_url": "https://s3.yourdomain.com/patches/app-uuid/1/aarch64-android.patch",
    "hash": "sha256-hash-of-patch"
  }
}
\end{lstlisting}
\textit{Note: If \texttt{patch\_number} in the request is equal to or greater than the active patch for that release/channel, return \texttt{"patch\_available": false}.}

\newpage

\section{Integration Guide}

To connect the official Shorebird tooling to your self-hosted backend, specific environment variables and configurations must be applied.

\subsection{CLI Configuration}
The Shorebird CLI relies on the \texttt{SHOREBIRD\_HOSTED\_URL} variable to determine the API endpoint.

\begin{lstlisting}[language=bash]
# Export the custom API URL
export SHOREBIRD_HOSTED_URL="https://api.yourdomain.com"

# Export the custom Artifact Proxy URL (for engine binaries)
export FLUTTER_STORAGE_BASE_URL="https://storage-proxy.yourdomain.com"

# Now run shorebird commands normally
shorebird init
shorebird release android
shorebird patch android
\end{lstlisting}

Alternatively, you can specify the \texttt{base\_url} directly in the \texttt{shorebird.yaml} file of your Flutter project:
\begin{lstlisting}[language=yaml]
app_id: your-app-id
base_url: https://api.yourdomain.com
\end{lstlisting}

\subsection{Artifact Proxy Setup}
The \texttt{artifact\_proxy} package included in the Shorebird repository must be deployed to intercept engine requests. You must configure its internal routing to point to your S3/MinIO buckets instead of Google Cloud Storage.
When the Flutter toolchain attempts to download \texttt{libflutter.so}, the proxy will serve the Shorebird-modified engine from your private bucket.

\subsection{Mobile App (Updater) Configuration}
When \texttt{shorebird init} generates the \texttt{shorebird.yaml} file, the Native Updater library reads this file at runtime. Because you defined \texttt{base\_url} in the YAML, the Rust-based updater will automatically direct its \texttt{/api/v1/patches/check} requests to your self-hosted backend. No native code modifications are required in the \texttt{updater} library itself.

\section{Conclusion}
By implementing this backend architecture, organizations can achieve a fully self-hosted Code Push environment. The design leverages the inherent flexibility of Shorebird's client-side tooling, requiring only standard HTTP API implementation and S3-compatible storage to achieve complete data sovereignty and operational control over the update lifecycle.

\end{document}
